\section{Discussion}\label{sec:discussion}

The hypotheses of Section~\ref{sec:intro} received different degrees of support. H1 was supported where Proposition~\ref{prop:failure} predicts it: $\varepsilon$-elimination lost the Pareto set in every replication of the cycle and in some replications of the chain, but not in the tie configuration, where the certificates that create the cycle appear in an order the procedure never reaches. The failure depends on the configuration and can pass a test suite unnoticed. H2 was supported for the proposed procedure, which made no false statement in any replication, but not for the comparison: repeated $t$ intervals erred in one configuration at a rate below $\alpha$, and confidence sequences without a correction did not err. What separates these methods in our studies is the guarantee, not the observed error rate. H3 was supported in direction only: the open-pair allocation saved evaluations in every cell, but the saving did not grow with $K$.

Without the union bound no error occurred because a false statement requires a sequence to fail in the direction that changes a certificate, when that certificate is decided, and only a few of the $Km$ sequences matter; the union bound protects against all failures and is conservative. Its cost is moderate, since widths grow with $\log(Km/\alpha)$: with $K=20$, $m=2$ and $\alpha=0.05$ it replaces $\log 20\approx3.0$ by $\log 800\approx6.7$. A sharper correction would have to use the correlation between objectives, which in Study~2 changed the number of evaluations by 10 to 15\%; rules for combining e-values \citep{vovk2021} are a possible route.

The budget of Theorem~\ref{thm:termination} is valid for every instance but is not a planning tool: for $K=10$, $m=2$ and $\varepsilon=0.05$ the Hoeffding sequence requires about $4.6\times10^6$ evaluations (Table~\ref{S-tab:horizon}), more than a thousand times the median used by the betting sequence in Study~2, for which we have no comparable bound. A pilot identification is a better guide, and, being anytime-valid, it is not wasted.

The setting has two limitations. First, the bounds of every objective must be known and must hold. For averages of indicators and other bounded measures they do; for costs or response times without a natural ceiling the analyst must justify a bound or transform the objective. In the inventory application a bound six times wider than the observed range cost little, but only because the betting sequence adapts to the variance. Second, evaluations of different alternatives are treated as independent, so common random numbers, which in ranking and selection often reduce the number of replications substantially \citep{hunter2019}, are not exploited. Anytime-valid sequences for paired differences exist \citep{seqbench}, and certificates built on differences would combine both ideas, at the price of $K(K-1)m/2$ sequences instead of $Km$.

The applications showed a property of the target that we had not anticipated: an alternative within $\varepsilon_i$ of the lower bound of objective $i$ is $\epsb$-Pareto-optimal whatever its other values, and five dominated inventory policies were accepted for this reason. The certified cover removes them, but its phase is more expensive and its last pairs, whose cover margins are near zero, did not resolve within the budget; Theorem~\ref{thm:termination} does not cover this phase. We regard $\Peps$ as a stopping criterion and the cover as the description of the frontier.

\subsection{Conclusions}\label{sec:conclusion}

Identifying the Pareto set from noisy simulations involves two decisions that the usual procedures merge: what to discard and when to stop. We separated them. Discarding uses only exact dominance certified by simultaneous confidence boxes, so it never loses a Pareto-optimal alternative; the tolerance enters only the decision to accept, and a certified cover summarizes the result. The statements hold with probability at least $1-\alpha$ simultaneously over time, under any allocation and at any stopping time. %CONCLNUM% The price is the need for known bounds and the conservativeness of the union bound; the gain is that every intermediate report, including the one produced when the budget runs out, is a valid statement about the alternatives.
